\par\Needspace{4\baselineskip}
\begin{center}
\SourceLeaf{046}\pdfbookmark[2]{30.}{article-030}30.
\end{center}
Da keine der Zahlen, die wir den logarithmischen und trigonometrischen Tafeln entnehmen, absolut genau sein kann und jede nur bis zu einem gewissen Grad angenähert ist, können alle mithilfe dieser Tafeln ausgeführten Rechnungen nur annähernd richtige Ergebnisse liefern. In den meisten Fällen bieten zwar die gewöhnlichen Tafeln, die bis zur siebten Dezimalstelle genau sind, also vom wahren Wert an der siebten Stelle weder nach oben noch nach unten um mehr als eine halbe Einheit abweichen, eine mehr als ausreichende Genauigkeit, sodass die unvermeidlichen Fehler ganz ohne Belang sind. Gleichwohl kann es in besonderen Fällen vorkommen, dass sich die Wirkung der Tafelfehler so sehr verstärkt, dass wir eine sonst optimale Methode ganz aufgeben und durch eine andere ersetzen müssen. Ein solcher Fall kann auch bei den Rechnungen auftreten, die wir bisher erläutert haben. Es entspricht daher unserem Vorhaben, hier einige Untersuchungen über den Grad der Genauigkeit anzustellen, den die gewöhnlichen Tafeln bei diesen Rechnungen zulassen. Obwohl wir hier dieses für das praktische Rechnen höchst wichtige Thema nicht erschöpfend behandeln können, wollen wir die Untersuchung so weit führen, dass sie für unseren Zweck ausreicht und von jedem Interessierten weiter ausgearbeitet und auf beliebige andere Rechenoperationen ausgedehnt werden kann.

\par\Needspace{4\baselineskip}
\begin{center}
\hypertarget{native-046-article031}{}\pdfbookmark[2]{31.}{article-031}31.
\end{center}
Jeder Logarithmus, Sinus, Tangens usw. (oder allgemeiner jede einer Tafel entnommene irrationale Größe) ist mit einem Fehler behaftet, der bis zu einer halben Einheit der letzten Stelle reichen kann. Diese Fehlergrenze bezeichnen wir mit $\omega$. Für die gewöhnlichen Tafeln ergibt sich damit der Wert $=0{,}000\,0000\,5$. Kann der Logarithmus usw. der Tafel nicht unmittelbar entnommen werden und muss man ihn durch Interpolation ermitteln, so kann der Fehler aus zwei Gründen noch etwas größer ausfallen. \emph{Erstens} setzt man für den proportionalen Teil, sooft er keine ganze Zahl ist (wenn man die letzte Stelle als Einheit nimmt), gewöhnlich die nächstgrößere oder nächstkleinere ganze Zahl ein. Wie man leicht sieht, kann sich der Fehler dadurch nahezu verdoppeln. Diese Vergrößerung des Fehlers lassen wir hier ganz außer Betracht, da wir dem proportionalen Teil ohne Weiteres noch eine oder zwei Dezimalstellen \SourceLeaf{047}anfügen können und unmittelbar ersichtlich ist, dass der interpolierte Logarithmus bei absolut genauem proportionalem Teil keinem größeren Fehler unterliegt als die unmittelbar in den Tafeln angegebenen Logarithmen, soweit sich deren Änderungen als gleichförmig auffassen lassen. Die \emph{zweite} Vergrößerung des Fehlers entsteht dadurch, dass diese Annahme streng genommen nicht zutrifft. Auch sie vernachlässigen wir, da der Einfluss der zweiten und höheren Differenzen in fast allen Fällen völlig unerheblich ist (besonders wenn man für die trigonometrischen Größen die hervorragenden Tafeln verwendet, deren Erstellung \textsc{Taylor} betreute) und sich ohne Mühe berücksichtigen lässt, wo er etwa etwas größer ausfällt. Wir setzen daher in allen Fällen den größten unvermeidlichen Tafelfehler $=\omega$, sofern das Argument (d. h. die Zahl, deren Logarithmus gesucht wird, oder der Winkel, dessen Sinus usw. gesucht wird) mit absoluter Genauigkeit bekannt ist. Ist dagegen das Argument selbst nur näherungsweise bekannt und nimmt man an, dass seinem größtmöglichen Fehler eine Änderung $\omega'$ des Logarithmus usw. entspricht (eine Änderung, die sich mithilfe von Differenzialen bestimmen lässt), so kann der größte Fehler des mithilfe der Tafeln berechneten Logarithmus bis zu $\omega+\omega'$ betragen.

Berechnet man umgekehrt mithilfe der Tafeln das Argument zu einem gegebenen Logarithmus, so ist sein größter Fehler gleich der Änderung des Arguments, die einer Änderung $\omega$ des Logarithmus entspricht, wenn dieser genau gegeben ist, oder einer Änderung $\omega+\omega'$ des Logarithmus, wenn dieser selbst mit einem Fehler von bis zu $\omega'$ behaftet sein kann. Es ist wohl kaum nötig zu bemerken, dass $\omega$ und $\omega'$ dasselbe Vorzeichen erhalten müssen.

Addiert man mehrere nur innerhalb bestimmter Grenzen genaue Größen, so ist der größte Fehler der Summe gleich der Summe der größten Einzelfehler, die sämtlich mit demselben Vorzeichen angesetzt werden. Daher ist auch bei der Subtraktion nur annähernd genauer Größen der größte Fehler der Differenz gleich der Summe der größten Einzelfehler. Bei der Multiplikation oder Division einer nicht absolut genauen Größe vergrößert oder verkleinert sich der größte Fehler im gleichen Verhältnis wie die Größe selbst.
